\documentclass[10pt,a4paper]{amsart}
\usepackage{amsmath,amssymb,amsthm,graphicx}
\usepackage{verbatim}
\newenvironment{reference}{\comment}{\endcomment}
%\newenvironment{reference}{}{}
\newenvironment{slogan}{\comment}{\endcomment}
\newenvironment{history}{\comment}{\endcomment}

% For commutative diagrams we use Xy-pic
\usepackage[all]{xy}

% We use 2cell for 2-commutative diagrams.
\xyoption{2cell}
% Theorem environments.
%
\theoremstyle{plain}
\newtheorem{theorem}[subsection]{Theorem}
\newtheorem{proposition}[subsection]{Proposition}
\newtheorem{lemma}[subsection]{Lemma}

\theoremstyle{definition}
\newtheorem{definition}[subsection]{Definition}
\newtheorem{example}[subsection]{Example}
\newtheorem{exercise}[subsection]{Exercise}
\newtheorem{situation}[subsection]{Situation}

\theoremstyle{remark}
\newtheorem{remark}[subsection]{Remark}
\newtheorem{remarks}[subsection]{Remarks}

\numberwithin{equation}{subsection}

% Macros
%
\def\lim{\mathop{\mathrm{lim}}\nolimits}
\def\colim{\mathop{\mathrm{colim}}\nolimits}
\def\Spec{\mathop{\mathrm{Spec}}}
\def\Hom{\mathop{\mathrm{Hom}}\nolimits}
\def\Ext{\mathop{\mathrm{Ext}}\nolimits}
\def\SheafHom{\mathop{\mathcal{H}\!\mathit{om}}\nolimits}
\def\SheafExt{\mathop{\mathcal{E}\!\mathit{xt}}\nolimits}
\def\Sch{\mathit{Sch}}
\def\Mor{\mathop{\mathrm{Mor}}\nolimits}
\def\Ob{\mathop{\mathrm{Ob}}\nolimits}
\def\Sh{\mathop{\mathit{Sh}}\nolimits}
\def\NL{\mathop{N\!L}\nolimits}
\def\CH{\mathop{\mathrm{CH}}\nolimits}
\def\proetale{{pro\text{-}\acute{e}tale}}
\def\etale{{\acute{e}tale}}
\def\QCoh{\mathit{QCoh}}
\def\Ker{\mathop{\mathrm{Ker}}}
\def\Im{\mathop{\mathrm{Im}}}
\def\Coker{\mathop{\mathrm{Coker}}}
\def\Coim{\mathop{\mathrm{Coim}}}

% Boxtimes
%
\DeclareMathSymbol{\boxtimes}{\mathbin}{AMSa}{"02}

%
% Macros for moduli stacks/spaces
%
\def\QCohstack{\mathcal{QC}\!\mathit{oh}}
\def\Cohstack{\mathcal{C}\!\mathit{oh}}
\def\Spacesstack{\mathcal{S}\!\mathit{paces}}
\def\Quotfunctor{\mathrm{Quot}}
\def\Hilbfunctor{\mathrm{Hilb}}
\def\Curvesstack{\mathcal{C}\!\mathit{urves}}
\def\Polarizedstack{\mathcal{P}\!\mathit{olarized}}
\def\Complexesstack{\mathcal{C}\!\mathit{omplexes}}
% \Pic is the operator that assigns to X its picard group, usage \Pic(X)
% \Picardstack_{X/B} denotes the Picard stack of X over B
% \Picardfunctor_{X/B} denotes the Picard functor of X over B
\def\Pic{\mathop{\mathrm{Pic}}\nolimits}
\def\Picardstack{\mathcal{P}\!\mathit{ic}}
\def\Picardfunctor{\mathrm{Pic}}
\def\Deformationcategory{\mathcal{D}\!\mathit{ef}}
\usepackage{cleveref}
\crefname{originallemma}{original lemma}{original lemmas}
\Crefname{originallemma}{Original lemma}{Original lemmas}
\makeatletter
\newlabel{algebra-lemma-minimal-prime-reduced-ring}{{00EU}{source}{}{sourcepointer-1}{}}
\newlabel{algebra-lemma-minimal-prime-reduced-ring@cref}{{[originallemma][0][]00EU}{[1][0][]source}{}{}{}}
\makeatother
\crefname{originallemma}{original lemma}{original lemmas}
\Crefname{originallemma}{Original lemma}{Original lemmas}
\makeatletter
\newlabel{decent-spaces-lemma-decent-generic-points}{{0ABV}{source}{}{sourcepointer-2}{}}
\newlabel{decent-spaces-lemma-decent-generic-points@cref}{{[originallemma][0][]0ABV}{[1][0][]source}{}{}{}}
\makeatother
\crefname{originalreference}{original reference}{original references}
\Crefname{originalreference}{Original reference}{Original references}
\makeatletter
\newlabel{lemma-proper-geometrically-reduced-global-sections}{{0BUG}{source}{}{sourcepointer-3}{}}
\newlabel{lemma-proper-geometrically-reduced-global-sections@cref}{{[originalreference][0][]0BUG}{[1][0][]source}{}{}{}}
\makeatother
\title{An invertible module on a proper integral algebraic space with sections of it and of its dual is trivial}
\author{The Stacks Project Authors}
\date{stacks-project:spaces-over-fields.tex:local-snapshot-20261012; GNU FDL 1.2 or later}
\begin{document}
\maketitle
\noindent\textit{Editorial note.} All delivered mathematical text is the author's own, copied from the pinned Stacks Project chapter file spaces-over-fields.tex; no statement or proof step was rewritten, reordered or removed.\par\medskip
\noindent\textit{Editorial note.} The counted proof is the author's complete argument: the argument that a nonzero section is nonzero at the generic point through trivialization on an affine etale neighbourhood, the same argument for the dual module, and the conclusion from the field property of the ring of global functions.\par\medskip
\noindent\textit{Editorial note.} The label lemma-characterize-trivial-pic-integral names the counted statement, and the proof environment immediately following it in the source is delivered verbatim as its proof.\par\medskip
\noindent\textit{Editorial note.} The counted statement is carried by the author's own lemma environment in the chapter source.\par\medskip
\noindent\textit{Editorial note.} Pre-open gate: python operations/agent-authoring/tools/precard\_gate.py --source spaces-over-fields.tex --label lemma-characterize-trivial-pic-integral --start 2724 --end 2749 --local-id 28063 returned RESULT=OK (L28063 free, occupancy FREE, public interval clear). Candidate taken from sources/stacks-spaces-over-fields/candidates-halfA.json.\par\medskip
\noindent\textit{Editorial note.} The statement and proof are copied from the pinned chapter file; the cited inputs are kept as external pointers with their canonical Stacks tags.\par\medskip
\noindent\textit{Allowed standard prerequisites (editorial).}\begin{itemize}
\item The ring of global functions on a proper geometrically reduced space is a field. Same chapter, the lemma showing that the global sections of the structure sheaf on a proper geometrically reduced algebraic space over a field form a field. Supplies the fact that the product of the two sections, viewed as a global function, lives in a field, hence is either zero or everywhere invertible.
\item Generic points of a decent space and their local rings. Decent Spaces, the lemma describing the local ring at a generic point of a reduced space as a field of dimension zero. Supplies the description of the local ring at the generic point, used to deduce that the product of the two nonzero sections is nonzero.
\item Nonvanishing of a regular function over a component's generic point. Algebra, the lemma that the local ring of a reduced ring at a minimal prime is a field, applied to the function corresponding to the section over an etale neighbourhood. Supplies the conclusion that a nonzero section is nonzero at the generic point of each irreducible component of an affine etale neighbourhood.
\end{itemize}\medskip
\ifdefined\hypertarget\hypertarget{sourcepointer-1}{}\fi
\noindent\textit{Source reference.} Original lemma 00EU: sources/stacks-chow/tags.txt, source lines 548--548. Prerequisite pointer invoked at proof line(s) 2744; it names a standard fact of a sibling chapter that the counted argument uses, and every step of the counted proof built on it is retained verbatim in the delivered body.\par\medskip
\ifdefined\hypertarget\hypertarget{sourcepointer-2}{}\fi
\noindent\textit{Source reference.} Original lemma 0ABV: sources/stacks-chow/tags.txt, source lines 13033--13033. Prerequisite pointer invoked at proof line(s) 2739; it names a standard fact of a sibling chapter that the counted argument uses, and every step of the counted proof built on it is retained verbatim in the delivered body.\par\medskip
\ifdefined\hypertarget\hypertarget{sourcepointer-3}{}\fi
\noindent\textit{Source reference.} Original reference 0BUG: sources/stacks-spaces-over-fields/pinned-original-source.tex, source lines 2646--2646. Prerequisite pointer invoked at proof line(s) 2746; it names a standard fact of the same chapter that the counted argument uses, and every step of the counted proof built on it is retained verbatim in the delivered body.\par\medskip
\begin{lemma}
\label{lemma-characterize-trivial-pic-integral}
Let $k$ be a field. Let $X$ be a proper integral algebraic space over $k$.
Let $\mathcal{L}$ be an invertible $\mathcal{O}_X$-module.
If $H^0(X, \mathcal{L})$ and $H^0(X, \mathcal{L}^{\otimes - 1})$
are both nonzero, then $\mathcal{L} \cong \mathcal{O}_X$.
\end{lemma}

\begin{proof}
Let $s \in H^0(X, \mathcal{L})$ and $t \in H^0(X, \mathcal{L}^{\otimes - 1})$
be nonzero sections. Let $x \in |X|$ be a point in the support of $s$.
Choose an affine \'etale neighbourhood $(U, u) \to (X, x)$ such that
$\mathcal{L}|_U \cong \mathcal{O}_U$. Then $s|_U$ corresponds to a nonzero
regular function on the reduced (because $X$ is reduced) scheme $U$ and hence
is nonvanishing in a generic point of an irreducible component of $U$. By
Decent Spaces, Lemma \ref{decent-spaces-lemma-decent-generic-points}
we conclude that the generic point $\eta$ of $|X|$ is in the support of $s$.
The same is true for $t$. Then of course $st$ must be nonzero because
the local ring of $X$ at $\eta$ is a field (by aforementioned lemma
the local ring has dimension zero, as $X$ is reduced the local ring is
reduced, and Algebra, Lemma \ref{algebra-lemma-minimal-prime-reduced-ring}).
However, we have seen that $K = H^0(X, \mathcal{O}_X)$
is a field in Lemma \ref{lemma-proper-geometrically-reduced-global-sections}.
Thus $st$ is everywhere nonzero and we see that
$s : \mathcal{O}_X \to \mathcal{L}$ is an isomorphism.
\end{proof}

\end{document}
